
\documentclass{article}
\usepackage{colortbl}
\usepackage{makecell}
\usepackage{multirow}
\usepackage{supertabular}

\begin{document}

\newcounter{utterance}

\centering \large Interaction Transcript for game `cladder', experiment `full\_v1.5\_default', episode 2927 with reflexion{-}Qwen3.5{-}9B.
\vspace{24pt}

{ \footnotesize  \setcounter{utterance}{1}
\setlength{\tabcolsep}{0pt}
\begin{supertabular}{c@{$\;$}|p{.15\linewidth}@{}p{.15\linewidth}p{.15\linewidth}p{.15\linewidth}p{.15\linewidth}p{.15\linewidth}}
    \# & \multicolumn{2}{c}{Player} && \multicolumn{2}{c}{Game Master} \\
    \hline

    \theutterance \stepcounter{utterance}  
    & & & \multicolumn{4}{p{0.6\linewidth}}{
        \cellcolor[rgb]{0.9,0.9,0.9}{
            \makecell[{{p{\linewidth}}}]{
                \texttt{\tiny{[P1$\langle$GM]}}
                \texttt{You have attempted this same task before and did not succeed. These are the plans you wrote after those attempts.} \\
\texttt{Trial \#1: New plan: I failed because I provided a detailed causal analysis and a "Yes" conclusion in the middle of my response, but I did not start the answer with the required "Yes" or "No" immediately as the first words. Additionally, I got cut off mid{-}sentence while attempting a self{-}correction, which likely caused the system to reject the incomplete output or fail to register the final correct answer. Next time, I will strictly adhere to the format by starting the very first token with "Yes" or "No", followed immediately by the reasoning on a new line, ensuring the response is complete and does not contain any introductory text or partial sentences.} \\
\\ 
\\ 
\texttt{You are an expert in causal inference. The following question is not a typical commonsense query, but rather a meticulously designed question created by a professor specializing in causal inference, intended to assess the students' mastery of the course content.} \\
\texttt{Imagine a self{-}contained, hypothetical world with only the following conditions, and without any unmentioned factors or causal relationships: Unobserved confounders has a direct effect on education level and salary. Proximity to a college has a direct effect on education level. Education level has a direct effect on salary. Unobserved confounders is unobserved. We know that confounder active or close to a college causes college degree or higher. confounder active and college degree or higher causes high salary. We observed the person lives close to a college and confounder active. Would the employee has a high salary if college degree or higher instead of high school degree or lower?} \\
\texttt{Start your answer with 'Yes' or 'No', followed by additional reasoning or evidence to support your explanation.} \\
            }
        }
    }
     \\ \\

    \theutterance \stepcounter{utterance}  
    & \multicolumn{4}{p{0.6\linewidth}}{
        \cellcolor[rgb]{0.561,0.737,0.561}{
            \makecell[{{p{\linewidth}}}]{
                \texttt{\tiny{[P1$\rangle$GM]}}
                \texttt{Yes} \\
\\ 
\texttt{In this hypothetical world, the causal structure involves unobserved confounders (\$U\$) affecting both education level (\$E\$) and salary (\$S\$), while proximity to a college (\$P\$) also directly affects education. The question asks for a counterfactual: given that an individual lives close to a college (\$P=1\$) and has an active confounder (\$U=1\$), would their salary be higher if they had a college degree (\$E=1\$) instead of a high school degree (\$E=0\$)?} \\
\\ 
\texttt{The prompt explicitly states the causal mechanisms: "confounder active and college degree or higher causes high salary." This implies that the combination of \$U=1\$ and \$E=1\$ is a sufficient condition for a high salary. Since the individual already possesses the active confounder (\$U=1\$), the only missing variable to trigger the "high salary" outcome is the education level. By changing the education level from high school or lower (\$E=0\$) to college degree or higher (\$E=1\$), the sufficient condition (\$U=1 \textbackslash{}land E=1\$) is met. Therefore, the employee would have a high salary under the counterfactual scenario. The direct effect of proximity to a college on education is irrelevant to the counterfactual comparison because the education level is being explicitly set to "college degree or higher" in the hypothetical condition, overriding whatever the proximity might have caused.} \\
            }
        }
    }
    & & \\ \\

    \theutterance \stepcounter{utterance}  
    & & & \multicolumn{2}{p{0.3\linewidth}}{
        \cellcolor[rgb]{0.9,0.9,0.9}{
            \makecell[{{p{\linewidth}}}]{
                \texttt{\tiny{[GM$|$GM]}}
                \texttt{Yes} \\
            }
        }
    }
    & & \\ \\

    \theutterance \stepcounter{utterance}  
    & & & \multicolumn{2}{p{0.3\linewidth}}{
        \cellcolor[rgb]{0.9,0.9,0.9}{
            \makecell[{{p{\linewidth}}}]{
                \texttt{\tiny{[GM$|$GM]}}
                \texttt{yes} \\
            }
        }
    }
    & & \\ \\

    \theutterance \stepcounter{utterance}  
    & & & \multicolumn{2}{p{0.3\linewidth}}{
        \cellcolor[rgb]{0.9,0.9,0.9}{
            \makecell[{{p{\linewidth}}}]{
                \texttt{\tiny{[GM$|$GM]}}
                \texttt{game\_result = WIN} \\
            }
        }
    }
    & & \\ \\

\end{supertabular}
}

\end{document}
